\chapter{Conclusion}\label{chapter:conclusion}

This thesis investigated the state-of-the-art in discrete variate generation for Zipfian distributions within benchmarking applications used in storage and database systems research.
 Four widely-used benchmarking tools were identified and analyzed: YCSB, FIO, Sysbench, and OLTP-Bench, each employing its own implementation of Zipfian variate generation.

In parallel, a survey of relevant literature in variate generation led to the identification of four promising methods: 
Rejection-Inversion Sampling, Rejection Sampling, Condensed Table Lookup, and Approximative Sampling. 
Of these, Rejection Sampling and Condensed Table Lookup were implemented, as the others are already in use within the studied tools.

The evaluation, which measured both performance and statistical accuracy, revealed that Sysbench and the 
Condensed Table Lookup method (both C-based) demonstrated strong performance. However, the latter lacked 
sufficient accuracy in the evaluated metrics, warranting further investigation before practical adoption. 
FIO was found to be both inaccurate and among the slowest implementations, making it unsuitable in its current form. 
Similarly, the Rejection Sampling method performed poorly in terms of both speed and accuracy. 

Based on these findings, this work recommends using YCSB and Sysbench 
for benchmarking tasks requiring Zipfian distributions. Additionally, 
updating FIO with an improved hashing function could make it a viable option as well.
